\documentclass[12pt,a4paper]{article}

% Кодування/мови — в такому порядку
\usepackage[utf8]{inputenc}
\usepackage[T2A]{fontenc}
\usepackage[ukrainian]{babel}
\usepackage{cmap} 

\usepackage{geometry}
\geometry{left=2cm,right=2cm,top=2cm,bottom=2cm}

\usepackage{graphicx}
\usepackage{array,multirow,booktabs,makecell}
\usepackage{caption}
\renewcommand{\thetable}{\arabic{table}}
\usepackage[table]{xcolor} % <-- це дає \rowcolor і \cellcolor
\captionsetup[table]{name=Таблиця}

\usepackage{amsmath}
\usepackage{pdflscape}
\usepackage{xcolor}

\newcommand{\g}{\cellcolor{gray!25}} % коротке позначення сірого

% Гіперпосилання — тримай наприкінці пакетів, з unicode
\usepackage[unicode]{hyperref}

\renewcommand\cellalign{tc}
\renewcommand\theadalign{cc}
\newcommand{\addh}{\rule{0pt}{3.8ex}} % ДОДАТКОВА ВИСОТА РЯДКА (лише де потрібно)
\renewcommand{\arraystretch}{1.25}

\begin{document}

    \begin{titlepage}

        % Картинка шапки (по центру, з відступом від верху)
        \begin{center}
\includegraphics[width=0.7\textwidth]{../../../TITLES/letterhead.png}\end{center}

        \begin{center}
            \large Національний технічний університет України\\
            «Київський політехнічний інститут імені Ігоря Сікорського»\\[1em]
            Факультет інформатики та обчислювальної техніки\\
            Кафедра обчислювальної техніки
        \end{center}

        \vfill

        \begin{center}
            \textbf{\LARGE Теорія електричних кіл та сигналів}\\[2em]
            \textbf{\Large Лабораторна робота №5}\\[0.5em]
            «Дослідження пасивного чотириполюсника змінного струму»
        \end{center}

        \vfill

        \begin{flushright}
            Виконав: студент 2 курсу ФІОТ, гр. ІО-41\\
            \textit{Давидчук А.~М.}\\[1em]
            Перевірив: \textit{Лободзинський В.~Ю.}
        \end{flushright}

        \vfill

        \begin{center}Київ --- 2025\end{center}

    \end{titlepage}

    \setlength{\parindent}{0pt}

    \textbf{\underline{Тема:}} «Дослідження пасивного чотириполюсника змінного струму».

    \vspace{1em}

    \textbf{\underline{Мета:}} Оволодіння методикою експериментального дослідження пасивного чотириполюсника
    змінного струму, визначення параметрів чотириполюсника та перевірка співвідношень,
    встановлених для цих параметрів.

    \begin{center} \textbf{\Large Практична частина} \end{center}

    \setlength{\parindent}{1.5em}

    Усі розрахунки було проведено за допомогою мови програмування Python (вихідний код додано разом із звітом).

    \begin{center} \textbf{\large Завдання №1 (офлайн вимірювання)} \end{center}

    Під час проведення офлайн заняття, ми скали таку електричну схему:



    %Моя група: ІО-41, мій порядковий номер у групі: 6. Звідси $N = 6$, $S = 1$.
    %Звідси загальна структурна схема електричного кола:

    \begin{figure}[h!]

        \renewcommand{\thefigure}{\arabic{figure}} % робимо "3.1", "3.2" і т.д.

        \centering
        % Підставляєте потрібний шлях та розмір зображення:
        \includegraphics[width=0.8\textwidth]{scheme0.png}
        % Підпис (зазвичай під малюнком):
        \caption{Структурна схема Т-подібного чотириполюсника}
        % Мітка для посилань у тексті (\ref{fig:...})
        \label{fig1:scheme0}

    \end{figure}

    Та проводили в лабораторії виміри сили струму, напруг, потужностей та зсуву фаз. Результат вимірів винесено в таблицю:


    \begin{table}[h!]
    \centering
    \resizebox{\textwidth}{!}{%
    %          1)Режим   2)label  3)U1  4)I1  5)P1  6)phi1  7)U2  8)I2  9)P2 10)phi2
    \begin{tabular}{|p{48mm}|c|c|c|c|c|c|c|c|c|}
    \hline
    \thead{Режим} & \multicolumn{1}{c|}{} &
    \thead{U$_1$, В} & \thead{I$_1$, A} & \thead{P$_1$, Вт} & \thead{$\varphi_1$,\\ град} &
    \thead{U$_2$, В} & \thead{I$_2$, A} & \thead{P$_2$, Вт} & \thead{$\varphi_2$,\\ град} \\
    \hline
    \multirow{2}{=}{\raggedright Живлення на вхідних\\ затискачах 1--1’}
    & н.х. & 12{,}07 & 0{,}0345 & 0{,}155 & $-68{,}14^\circ$ & \g & \g & \g & \g \\
    \cline{2-10}
    & к.з. & 12{,}30 & 0{,}0390 & 0{,}12 & $-75{,}51^\circ$ & \g & \g & \g & \g \\
    \hline
    \multirow{2}{=}{\raggedright Живлення\\ на вихідних\\ затискачах 2--2’}
    & \g \addh н.х. & \g & \g & \g & \g & 10{,}80 \addh & 0{,}049 \addh & 0{,}48 \addh & $24{,}90^\circ$ \addh \\
    \cline{2-10}
    & \g \addh к.з. & \g & \g & \g & \g & 10{,}60 \addh & 0{,}054 \addh & 0{,}56 \addh & $11{,}95^\circ$ \addh \\
    \hline
    \end{tabular}%
    }
    \caption{\centering Результати вимірювань фіз. величин в різних режимах роботи чотириполюсника}
    \end{table}

    А також було окремо виміряно параметри елементів чотириполюсника:

    \begin{table}[h!]
    \centering
    \setlength{\tabcolsep}{16pt}      % горизонтальні відступи в клітинках
    \renewcommand{\arraystretch}{1.6} % висота рядків

    \begin{tabular}{|c|c|c|c|c|}
    \hline
    R, Ом & RL, Ом & L, Гн & RC, Ом & C, мкФ \\
    \hline
       100         &        79{,}80          &        0{,}362         & \text{---}          &      9{,}43           \\
    \hline
    \end{tabular}
    \caption{}
    \end{table}

    \newpage

    За результатами дослідів обчислюємо комплексні значення вхідних опорів та перевіряємо справедливість співвідношення:

    \[
    \boxed{
    \frac{Z_{1HX}}{Z_\text{1КЗ}} = \frac{Z_{2HX}}{Z_\text{2КЗ}}
    }.
    \]

    де 

    \[
    Z_{1HX} = \frac{U_{1HX}}{I_{1HX}} = \frac{U_{1HX}}{I_{1HX}} \cdot 
    \frac{e^{i\psi_{U1HX}}}{e^{i\psi_{I1HX}}} = Z_{1HX} e^{i\varphi_{1HX}}
    = 349{,}855e^{-i68{,}14^\circ}
    \]

    \[
    Z_{\text{1КЗ}} = \frac{U_{\text{1КЗ}}}{I_{\text{1КЗ}}} = \frac{U_{\text{1КЗ}}}{I_{\text{1КЗ}}} \cdot 
    \frac{e^{i\psi_{U1\text{КЗ}}}}{e^{i\psi_{I1\text{КЗ}}}} = Z_{\text{1КЗ}} e^{i\varphi_{\text{1КЗ}}}
    = 315{,}385e^{-i75.51^\circ}
    \]

    \[
    Z_{2HX} = \frac{U_{2HX}}{I_{2HX}} = Z_{2HX} e^{i\varphi_{2HX}}
    = 220{,}408e^{i24{,}90^\circ}
    \]

    \[
    Z_{\text{2КЗ}} = \frac{U_{\text{2КЗ}}}{I_{\text{2КЗ}}} = Z_{\text{2КЗ}} e^{i\varphi_{\text{2КЗ}}}
    = 196{,}296e^{i11{,}95^\circ}
    \]

    І перевіряємо співвідношення:

    \[
    \frac{Z_{1HX}}{Z_{\text{1KЗ}}} = 1{,}109e^{i7{,}37^\circ} \text{ та }
    \frac{Z_{2HX}}{Z_{\text{2KЗ}}} = 1{,}123e^{i12{,}95^\circ}
    \]

    Що межах допустимої похибки обчислення та, авжеж, вимірювання.

    Тепер розрахуємо коефіцієнти чотириполюсника на основі попередніх даних за такими формулами:

    \[
    D = \sqrt{\frac{Z_{2HX}}{\,Z_{1HX}-Z_{1\text{KЗ}}\,}}
    = 1{,}848+0{,}776i
    \]

    \[
    C = \frac{D}{Z_{2HX}}
    = 0{,}009-0{,}000i
    \]

    \[
    A = Z_{1HX} \cdot C
    = 1{,}074-2{,}994i
    \]

    \[
    B = D \cdot Z_{\text{1KЗ}} = 
    382{,}690-503{,}049i
    \]

    Перевіряємо виконання умови $AD - BC = 1$:

    \[
    4{,}307-4{,}700i - 3{,}307+4{,}700i = 1{,}000+0{,}000i
    \]

    Тепер на основі коефіцієнтів чотириполюсника, визначимо його параметри, та порівняємо із виміряниими раніше (у таблиці №2):

    \[
    Y_0 = C = 0{,}009-0{,}000i \,\text{См} = 109{,}902+4{,}084i \,\text{Ом};
    \]
    \[
    Z_1 = \frac{A - 1}{C} = 20{,}363-328{,}783i \,\text{Ом}; \qquad 
    Z_2 = \frac{D - 1}{C} = 90{,}018+88{,}716i \,\text{Ом}.
    \]

    Порівнюючи результат обчислення імпедансів отримуємо, що маємо похибку між виміряними та обчисленими значеннями через коефіцієнти чотириполюсника.
    Це може бути пов'язано із похибка вимірювальних приладів, похибка ціни поділки, а також похибка округлення при обчисленнях. 

    \newpage

    \begin{center} \textbf{\large Завдання №2 (Завдання на навчально-дослідну роботу студента)} \end{center}

    Обчислити параметри елементів заданого електричного кола (рис. 2):

    \begin{figure}[h!]

        \renewcommand{\thefigure}{\arabic{figure}} % робимо "3.1", "3.2" і т.д.

        \centering
        % Підставляєте потрібний шлях та розмір зображення:
        \includegraphics[width=0.8\textwidth]{scheme1.png}
        % Підпис (зазвичай під малюнком):
        \caption{Структурна схема П-подібного чотириполюсника}
        % Мітка для посилань у тексті (\ref{fig:...})

    \end{figure}

    за такими формулами:

    \[
    Z_1 = (N \cdot S) - i(10 \cdot N) \cdot k,\ \text{Ом};
    \]
    \[
    Z_2 = (N + S) + i(10 \cdot S) \cdot k,\ \text{Ом};
    \]
    \[
    Z_0 = 10 \cdot S + N,\ \text{Ом};
    \]
    \[
    U_2 = 4 + S,\ \text{В};
    \]
    \[
    f = 10 \cdot S + 10 \cdot N,\ \text{Гц}.
    \]

    \noindent
    Для мого варіанту (у мене $S = 1$, $N = 6$, $k = 1$):

    $Z_1 = 6 - 60i$ Ом, $Z_2 = 7 + 10i$ Ом, $Z_0 = 16$ Ом, $U_1 = 20$ В, $U_2 = 5$ В, $f = 70$ Гц.

    \vspace{1em}

    \noindent
    Тепер розрахуємо A-параметри для цього П-подібного чотириполюсника:

    \[
    A = 1 + \frac{Z_0}{Z_2} = 1{,}752-1{,}074i
    \]

    \[
    B = Z_0 = 16;
    \]

    \[
    C = \frac{1}{Z_1} + \frac{1}{Z_2} + \frac{Z_0}{Z_1 Z_2} = 0{,}068-0{,}040i;
    \]

    \[
    D = 1 + \frac{Z_0}{Z_1} = 1{,}026+0{,}264i.
    \]

    Перевірка умови $AD - BC = 1$:

    \[
    AD - BC = 2{,}081-0{,}640j - 1{,}081-0{,}640j = 1{,}000-0{,}000j
    \]

    І так як відомі напруга на вході та напруга на виході, то скористаємось рівнянням
    пасивного чотириполюсника для Y форми:

    \[
    \dot{I}_1 = Y_{11}\,\dot{U}_1 + Y_{12}\,\dot{U}_2,
    \qquad
    \dot{I}_2' = Y_{21}\,\dot{U}_1 + Y_{22}\,\dot{U}_2
    \]

    \[
    Y_{11}=\frac{D}{B}, \qquad
    Y_{12}=-\frac{1}{B}, \qquad
    Y_{22}=\frac{A}{B}, \qquad
    Y_{21}=Y_{12}
    \]

    \[
    I_1 = Y_{11}U_1 + Y_{12}U_2 = 0{,}971+0{,}330i,
    \qquad
    |I_1| = 1{,}025,
    \qquad
    \arg(I_1) = 18{,}78^{\circ}
    \]

    \[
    I_2 = -\bigl(Y_{21}U_1 + Y_{22}U_2\bigr) = 0{,}703+0{,}336i,
    \qquad
    |I_2| = 0{,}779,
    \qquad
    \arg(I_2) = 25{,}53^{\circ}
    \]

    Навантаження при якому на зажимах 2-2’ буде напруга $U_2=5$:

    \[
    Z_H = \frac{U_2}{I_2} = 5{,}795-2{,}768j \, \text{Ом}
    \]

    Залишилось визначити параметри елементів чотириполюсника, побудувати його еквівалентну схему у Multisim та порівняти виміряні дані, із обчисленими.
    За варіантом, ми працюємо зі змінною напругою із частотою $f = 70$ Гц. Звідси $\omega = 2\pi f = 140\pi$. Звідси:

    \[
    C_1 = -\frac{1}{\omega\Im(Z_1)} = 37{,}894\,\text{мкФ}
    \]
    \[
    L_2 = \frac{\Im(Z_2)}{\omega} = 22{,}736 \,\text{мГн}
    \]
    \[
    C_H = -\frac{1}{\omega\Im(Z_H)} = 821{,}531 \,\text{мкФ}
    \]
    \[
    R_0 = 16 \,\text{Ом}\quad R_1 = 6 \,\text{Ом}\quad R_2 = 7 \,\text{Ом}\quad R_H = 5{,}795 \,\text{Ом}\quad
    \]

    \begin{figure}[h!]

        \renewcommand{\thefigure}{\arabic{figure}} % робимо "3.1", "3.2" і т.д.

        \centering
        \includegraphics[width=1.0\textwidth]{scheme.pdf}
        \caption{Ел. схема П-подібного чотириполюсника у Multisim}

    \end{figure}

    Як видно, із схеми у Multisim, виміряні значення струмів та напруг співпадають із обчисленими раніше.

    \vspace{1em}

    \noindent
    \textbf{\underline{Висновок:}} У ході виконання лабораторної роботи було досліджено пасивний чотириполюсник змінного струму у двох постановках: 
    експериментальне визначення параметрів Т-подібного чотириполюсника та аналітичний розрахунок параметрів 
    П-подібного чотириполюсника за заданими умовами. Під час офлайн-вимірювань були отримані значення струмів, напруг, 
    активних потужностей та кутів зсуву фаз у режимах холостого ходу та короткого замикання. На їх основі визначено 
    комплексні вхідні опори $Z_{1HX}$, $Z_{1\text{КЗ}}$, $Z_{2HX}$, $Z_{2\text{КЗ}}$ та перевірено співвідношення 
    $\frac{Z_{1HX}}{Z_{1\text{КЗ}}} = \frac{Z_{2HX}}{Z_{2\text{КЗ}}}$, яке у межах експериментальної похибки виконується, що підтверджує 
    узгодженість виміряних даних.

    На основі отриманих величин були визначені A-параметри Т-подібного чотириполюсника, а також знайдено 
    еквівалентні параметри $Y_0$, $Z_1$, $Z_2$. Порівняння цих параметрів із попередньо виміряними показало наявність 
    невеликих розбіжностей, зумовлених похибками вимірювальних приладів, неточністю фазових вимірів і округленням 
    комплексних величин.

    У другій частині роботи були обчислені параметри П-подібного чотириполюсника за варіантом. Знайдено його 
    A-параметри, вхідні струми та навантаження, а також побудовано еквівалентну схему, яку перевірено у середовищі 
    Multisim. Отримані експериментальні та симуляційні дані добре узгоджуються між собою і підтверджують умову 
    $AD - BC = 1$, що є критерієм фізичної реалізовності пасивного чотириполюсника.

    Таким чином, у роботі було закріплено навички визначення параметрів пасивних чотириполюсників, роботи з 
    комплексними величинами, переходу між різними типами параметрів (A-, Y-, Z-параметри), а також побудови 
    еквівалентних схем. Усі отримані результати підтверджують теоретичні закономірності та демонструють 
узгодженість експериментальних і розрахункових даних.

\end{document}